%% ===============================================================
\section{Experiments}
\label{sec:experiments}
The v2 evaluation has not been run.
Every new experimental number below is \ph{}.
Numbers copied from the root-cause report dated 13~October~2026 are marked as such.
They were computed on the previous pipeline and are not v2 results.
Exploratory outputs of \texttt{offline\_v2} are from that old pipeline, contain known bugs, and are not the main table.
The section is a measurement study, a main table, ablations, and analysis.
No cell has been filled from a v2 run.

\subsection{Experimental Setup}
\label{sec:setup}
\textit{Splits.}
WebQA calibration is at least 500 questions and the development split is 256 questions.
Both are drawn from the training split, they are disjoint, and neither is drawn from the validation split.
The test split is all 4{,}966 validation questions.
The previous statement that calibration came from the validation split, while the reported set was the entire validation split, is withdrawn.
OK-VQA uses 500 calibration questions, 256 development questions, and the remainder of the validation split as the test set.
The three OK-VQA splits are pairwise disjoint.
An empty test split is an error.
When the validation split has 5{,}046 questions, the remainder is 4{,}290.
MMQA and video are generalization settings whose sizes follow the available assets.
The current command interface does not include them, and their scores are \ph{}.

\textit{Settings and metrics.}
The primary WebQA setting is the distractor pool (\texttt{webqa\_full}), in which each question brings its own candidates.
Open retrieval (\texttt{webqa\_open\_*}; the current interface generates \texttt{webqa\_open\_img} at stage~s2) is the second setting.
OK-VQA is scored with official VQA accuracy only.
WebQA's primary metric is acc$_{\mathrm{approx}}$.
Normalized BART and QA$_{\mathrm{approx}}$ are reported as well, on the full set and on the image and text subsets.
The image subset has 2{,}511 questions and the text subset has 2{,}455.
Loose lexical WebQA scores are not placed next to acc$_{\mathrm{approx}}$.

\textit{Actions and the strong baseline.}
The decoded actions are \texttt{direct}, \texttt{static\_text}, \texttt{static\_image}, \texttt{static\_mixed}, and \texttt{rerank\_static}.
Adaptive comparators are an NTE trigger, SeaKR, the previous $\Uomni$ gate (v1), and UniversalRAG routing.
They are not eligible to be the strong baseline.
The strong baseline is the single action with the best calibration-set score.
Primary relative gain uses that score as the denominator.
The tables also report gain versus \texttt{static\_mixed}.
The oracle is the per-question best action in the same set.
Evidence Attention Bias is not a main-method arm.
All actions in a run share one prompt, one budget, and one candidate pool.
The prompt is frozen on the development split.
The test split is scored once after \texttt{docs/dev\_decisions.json} is frozen.
Grids are not selected on the test split.
The operating point previously taken from a test-set grid (Table~20 in the audited draft) is withdrawn.

\textit{Statistics.}
Uncertainty is a paired bootstrap 95\% confidence interval over questions.
Multiple comparisons use Holm correction.
Decoding is greedy, so a configuration is deterministic given its inputs.
This draft does not report a seed standard deviation under greedy decoding.
The fraction of questions on which the chosen action scores below the strong baseline is reported with the risk-control claim.
A score below \texttt{direct} is not described as better than no retrieval.

\textit{Development gates.}
Two preregistered gates are criteria, not results, and this draft does not claim that either has passed.
After the pipeline fixes, the development-set per-question oracle is compared with the strongest single action against a relative gap of 8\%.
A selector is eligible for the frozen test run when, on the development split, its relative gain over that baseline is at least 5\%, the lower end of its paired interval is positive, and its degradation rate is at most $\delta$.
The test split is not used to make those choices.

\textit{Figures withdrawn by the root-cause report.}
The following numbers are copied from the root-cause report of 13~October~2026.
They are not v2 results, and they are not written into Tables~\ref{tab:webqa}--\ref{tab:cost}.
Relative gains of +33\% on WebQA and +15\% on OK-VQA used historical cached baselines: a cached static score of 17.9 against 23.4 for static retrieval on the same 512-question batch, a gap of 5.5 points, under different prompts, budgets, and retrieval scores.
The same report withdraws a trained controller's +8.43\% and a +14.4\% chosen on a 32-question development set.
On a 256-question OK-VQA development slice, official VQA accuracy was 61.5 for direct answering, 36.3 for static retrieval, 48.1 for the gated static arm, and 36.3 for EAB.
Direct answering won 82 questions and lost 10 against static retrieval.
The per-question maximum of direct and static on that slice was 64.4, which is 2.9 points above direct answering (about +4.7\% relative).
The gated arm and EAB are both below direct answering, so the report does not call that slice an improvement over no retrieval.
On 512 WebQA questions, static retrieval scored 59.1 acc$_{\mathrm{approx}}$, 27.4 normalized BART, and 23.4 QA.
Image questions scored 63.2 acc$_{\mathrm{approx}}$ and 13.8 QA.
Text questions scored 55.0 acc$_{\mathrm{approx}}$ and 33.0 QA.
Over 2{,}560 items, a loose WebQA scorer and official acc$_{\mathrm{approx}}$ differed by 10.5 points on average (0.699 versus 0.594).
On one static OK-VQA batch, a loose scorer gave 65.8 where official VQA accuracy was about 47.9.
A greedy-decoding figure of $\pm 0.2$ is smaller than the binomial standard error of about 0.66 points on 5{,}046 OK-VQA questions, and it is not a confidence interval.
The v2 measurement that replaces these comparisons is the placeholder tables below.

\subsection{Measurement Study}
\label{sec:measure}
The measurement reserves three quantities.
The gap between the best single action and the per-question oracle is \ph{}.
Image questions and text questions are separate columns in Tables~\ref{tab:webqa} and~\ref{tab:okvqa}, and every column is \ph{}.
The size of the normalized BART factor, reported beside acc$_{\mathrm{approx}}$, is \ph{}.
A 50-question image-evidence diagnostic is specified, and its conclusion is \ph{}.
The audit's 256-question and 512-question figures above are the motive for these tables.
They are not substitutes for the cells.

\begin{table*}[t]
\caption{WebQA distractor test split (\texttt{webqa\_full}): placeholder measurement.
acc$_{\mathrm{approx}}$ is primary.
Normalized BART and QA$_{\mathrm{approx}}$ are secondary.
Image and text are subsets of the same run.
The strong baseline is whichever single action scores highest on the calibration split; its identity is frozen before test and is not chosen on this table.
The oracle is the per-question best action.
Every entry is \ph{}.
The v2 run has not been executed.}
\label{tab:webqa}
\scriptsize
\setlength{\tabcolsep}{3.5pt}
\begin{tabular}{@{}lccccccccc@{}}
\toprule
Action & \multicolumn{3}{c}{acc$_{\mathrm{approx}}$} & \multicolumn{3}{c}{norm.\ BART} & \multicolumn{3}{c}{QA$_{\mathrm{approx}}$} \\
\cmidrule(lr){2-4}\cmidrule(lr){5-7}\cmidrule(l){8-10}
 & All & Image & Text & All & Image & Text & All & Image & Text \\
\midrule
\texttt{direct} & \ph{} & \ph{} & \ph{} & \ph{} & \ph{} & \ph{} & \ph{} & \ph{} & \ph{} \\
\texttt{static\_text} & \ph{} & \ph{} & \ph{} & \ph{} & \ph{} & \ph{} & \ph{} & \ph{} & \ph{} \\
\texttt{static\_image} & \ph{} & \ph{} & \ph{} & \ph{} & \ph{} & \ph{} & \ph{} & \ph{} & \ph{} \\
\texttt{static\_mixed} & \ph{} & \ph{} & \ph{} & \ph{} & \ph{} & \ph{} & \ph{} & \ph{} & \ph{} \\
\texttt{rerank\_static} & \ph{} & \ph{} & \ph{} & \ph{} & \ph{} & \ph{} & \ph{} & \ph{} & \ph{} \\
Strong baseline & \ph{} & \ph{} & \ph{} & \ph{} & \ph{} & \ph{} & \ph{} & \ph{} & \ph{} \\
Oracle & \ph{} & \ph{} & \ph{} & \ph{} & \ph{} & \ph{} & \ph{} & \ph{} & \ph{} \\
\bottomrule
\end{tabular}
\end{table*}

\begin{table}[t]
\caption{OK-VQA test split: official VQA accuracy only.
Placeholder measurement.
The same action definitions as Table~\ref{tab:webqa}.
Every entry is \ph{}.}
\label{tab:okvqa}
\small
\begin{tabular}{@{}lc@{}}
\toprule
Action & VQA accuracy \\
\midrule
\texttt{direct} & \ph{} \\
\texttt{static\_text} & \ph{} \\
\texttt{static\_image} & \ph{} \\
\texttt{static\_mixed} & \ph{} \\
\texttt{rerank\_static} & \ph{} \\
Strong baseline & \ph{} \\
Oracle & \ph{} \\
\bottomrule
\end{tabular}
\end{table}

\subsection{Main Results}
\label{sec:main}
Table~\ref{tab:main} is the main comparison.
Relative gain versus the strong baseline is the primary column.
Gain versus \texttt{static\_mixed} is secondary.
Each gain cell is a point estimate and a paired bootstrap 95\% interval; both are \ph{}, so the cell is written \ph{} [\ph{},\,\ph{}].
Degradation is the fraction of test questions on which the chosen action scores below the strong baseline.
No row is marked significant, because no test has been run.
A filled row that sits below \texttt{direct} will not be described as better than no retrieval.

\begin{table*}[t]
\caption{Main results on the frozen test splits.
WebQA uses acc$_{\mathrm{approx}}$ in the distractor setting.
OK-VQA uses official VQA accuracy.
``Versus strong'' uses the calibration-set best single action as the denominator.
``Versus static'' uses \texttt{static\_mixed}.
Brackets are paired bootstrap 95\% intervals.
Degr.\ is the degradation rate against the strong baseline.
All entries are \ph{}.
EAB is omitted.}
\label{tab:main}
\scriptsize
\setlength{\tabcolsep}{2.4pt}
\begin{tabular}{@{}lccccccc@{}}
\toprule
Method & \multicolumn{3}{c}{WebQA acc$_{\mathrm{approx}}$} & \multicolumn{3}{c}{OK-VQA} & Degr. \\
\cmidrule(lr){2-4}\cmidrule(lr){5-7}
 & Score & vs.\ strong & vs.\ static & Score & vs.\ strong & vs.\ static & \\
\midrule
\texttt{direct} & \ph{} & \ph{} [\ph{},\,\ph{}] & \ph{} [\ph{},\,\ph{}] & \ph{} & \ph{} [\ph{},\,\ph{}] & \ph{} [\ph{},\,\ph{}] & \ph{} \\
\texttt{static\_text} & \ph{} & \ph{} [\ph{},\,\ph{}] & \ph{} [\ph{},\,\ph{}] & \ph{} & \ph{} [\ph{},\,\ph{}] & \ph{} [\ph{},\,\ph{}] & \ph{} \\
\texttt{static\_image} & \ph{} & \ph{} [\ph{},\,\ph{}] & \ph{} [\ph{},\,\ph{}] & \ph{} & \ph{} [\ph{},\,\ph{}] & \ph{} [\ph{},\,\ph{}] & \ph{} \\
\texttt{static\_mixed} & \ph{} & \ph{} [\ph{},\,\ph{}] & \ph{} [\ph{},\,\ph{}] & \ph{} & \ph{} [\ph{},\,\ph{}] & \ph{} [\ph{},\,\ph{}] & \ph{} \\
\texttt{rerank\_static} & \ph{} & \ph{} [\ph{},\,\ph{}] & \ph{} [\ph{},\,\ph{}] & \ph{} & \ph{} [\ph{},\,\ph{}] & \ph{} [\ph{},\,\ph{}] & \ph{} \\
NTE trigger & \ph{} & \ph{} [\ph{},\,\ph{}] & \ph{} [\ph{},\,\ph{}] & \ph{} & \ph{} [\ph{},\,\ph{}] & \ph{} [\ph{},\,\ph{}] & \ph{} \\
SeaKR & \ph{} & \ph{} [\ph{},\,\ph{}] & \ph{} [\ph{},\,\ph{}] & \ph{} & \ph{} [\ph{},\,\ph{}] & \ph{} [\ph{},\,\ph{}] & \ph{} \\
v1 $\Uomni$ gate & \ph{} & \ph{} [\ph{},\,\ph{}] & \ph{} [\ph{},\,\ph{}] & \ph{} & \ph{} [\ph{},\,\ph{}] & \ph{} [\ph{},\,\ph{}] & \ph{} \\
UniversalRAG & \ph{} & \ph{} [\ph{},\,\ph{}] & \ph{} [\ph{},\,\ph{}] & \ph{} & \ph{} [\ph{},\,\ph{}] & \ph{} [\ph{},\,\ph{}] & \ph{} \\
v2-Rule & \ph{} & \ph{} [\ph{},\,\ph{}] & \ph{} [\ph{},\,\ph{}] & \ph{} & \ph{} [\ph{},\,\ph{}] & \ph{} [\ph{},\,\ph{}] & \ph{} \\
v2-Cal & \ph{} & \ph{} [\ph{},\,\ph{}] & \ph{} [\ph{},\,\ph{}] & \ph{} & \ph{} [\ph{},\,\ph{}] & \ph{} [\ph{},\,\ph{}] & \ph{} \\
Oracle & \ph{} & \ph{} [\ph{},\,\ph{}] & \ph{} [\ph{},\,\ph{}] & \ph{} & \ph{} [\ph{},\,\ph{}] & \ph{} [\ph{},\,\ph{}] & \ph{} \\
\bottomrule
\end{tabular}
\end{table*}

\subsection{Ablations}
\label{sec:ablation}
Table~\ref{tab:ablation} removes one v2-Cal feature at a time, compares v2-Rule with v2-Cal, and leaves a row for the preregistered level $\alpha=\delta=0.10$.
The $\delta$ row is a development-split description.
It is not a point selected on the test set.
Refitting uses the calibration split only.
Every entry is \ph{}.

\begin{table}[t]
\caption{Ablations on the frozen test metrics after calibration-split refits.
WebQA is acc$_{\mathrm{approx}}$.
OK-VQA is official VQA accuracy.
Utility is Equation~\eqref{eq:loss} on the calibration split and is not an accuracy.
Every entry is \ph{}.
The $\delta$ row is not used to pick the test operating point.}
\label{tab:ablation}
\scriptsize
\setlength{\tabcolsep}{3pt}
\begin{tabular}{@{}lcccc@{}}
\toprule
Variant & WebQA & OK-VQA & Degr. & Utility \\
\midrule
v2-Cal, six features & \ph{} & \ph{} & \ph{} & \ph{} \\
Without normalized entropy & \ph{} & \ph{} & \ph{} & \ph{} \\
Without logit margin & \ph{} & \ph{} & \ph{} & \ph{} \\
Without sequence log-prob. & \ph{} & \ph{} & \ph{} & \ph{} \\
Without length & \ph{} & \ph{} & \ph{} & \ph{} \\
Without KL vs.\ direct & \ph{} & \ph{} & \ph{} & \ph{} \\
Without retriever gap & \ph{} & \ph{} & \ph{} & \ph{} \\
v2-Rule & \ph{} & \ph{} & \ph{} & \ph{} \\
Development $\delta=0.10$ & \ph{} & \ph{} & \ph{} & \ph{} \\
\bottomrule
\end{tabular}
\end{table}

\subsection{Open Retrieval, Generalization, and Cost}
\label{sec:extra}
Table~\ref{tab:open} is the open-retrieval image setting.
The current interface builds it with stage~s2, not with a frozen stage-s4 test, so the table is not a test result.
Generalization to MMQA and to video is \ph{} and has no command in that interface.
Table~\ref{tab:cost} reserves mean generator calls, input tokens, and latency.
Planning estimates of GPU time are not entered here.

\begin{table}[t]
\caption{Open-retrieval WebQA image setting (\texttt{webqa\_open\_img}), stage~s2 interface only.
Not a frozen test.
Primary metric acc$_{\mathrm{approx}}$, plus normalized BART and QA$_{\mathrm{approx}}$.
Every entry is \ph{}.}
\label{tab:open}
\small
\begin{tabular}{@{}lccc@{}}
\toprule
Action & acc$_{\mathrm{approx}}$ & norm.\ BART & QA$_{\mathrm{approx}}$ \\
\midrule
\texttt{direct} & \ph{} & \ph{} & \ph{} \\
\texttt{static\_text} & \ph{} & \ph{} & \ph{} \\
\texttt{static\_image} & \ph{} & \ph{} & \ph{} \\
\texttt{static\_mixed} & \ph{} & \ph{} & \ph{} \\
\texttt{rerank\_static} & \ph{} & \ph{} & \ph{} \\
\bottomrule
\end{tabular}
\end{table}

\begin{table}[t]
\caption{Cost placeholders on the stage-s4 runs.
Calls, input tokens, and latency are \ph{}.
The cheap variant decodes evidence actions only when the direct signal is below its threshold.}
\label{tab:cost}
\small
\begin{tabular}{@{}lccc@{}}
\toprule
System & Calls & Input tokens & Latency \\
\midrule
\texttt{direct} only & \ph{} & \ph{} & \ph{} \\
Always \texttt{static\_mixed} & \ph{} & \ph{} & \ph{} \\
SeaKR & \ph{} & \ph{} & \ph{} \\
v2, all actions & \ph{} & \ph{} & \ph{} \\
v2, cheap variant & \ph{} & \ph{} & \ph{} \\
\bottomrule
\end{tabular}
\end{table}

\subsection{Attention steering moves attention, not answers}
\label{sec:steer}
Every number in this subsection is copied from the root-cause report dated 13~October~2026.
The measurements come from the previous pipeline.
They are not v2 results, and they are not taken from \texttt{offline\_v2}.

On a 64-question probe, the L1 distance of attention toward the supporting block moved from 0.77 to 0.26, and the answer flipped on 2 of 64 questions.
On the 38 questions in that probe where the gate was on, the median gap between the top-1 and top-2 first-token logits was about 3.4--3.5, and the median largest logit displacement caused by the bias was about 2.0.
On 29 of those 38 questions the gap exceeded the displacement.
The two flips had gaps of 0 and 0.75.
First-token predictive entropy was about 0.01--0.02, and on 65\% of questions $h_{\mathrm{sem}}<0.01$.
Before the bias, about 59\% of the attention already lay on the supporting evidence.
The bias raised that absolute mass by 0.011 on average.
The largest bias was about 1.2 nats (median of the per-question maximum 0.78, 90th percentile 1.07).
Among 306 gated-on questions, the median attention odds between the best and worst chunk were 1.185, against the propositional upper bound $e^{1.2}=3.32$ at $U=1$ and a full relevance gap.
The same gap is why a uniform span bias in the style of PASTA, or a positional recalibration in the style of Found-in-the-Middle, has little room to change an answer the first token has already determined.
This draft does not report a new PASTA or Found-in-the-Middle run.

On 512 questions decoded under one protocol on CUDA, the first token changed on 14 questions.
EAB-P versus static retrieval changed official QA by +1.30\%, with a 95\% interval of $[-0.11,+0.87]$ points, which includes zero.
The same comparison on MPS was +1.20\% and also includes zero.
Of a +0.304 point QA difference, two questions contributed +0.195 and +0.122, and the other 510 questions summed to $-0.013$.
Replacing CUDA with MPS and adding no bias changed 62 of 512 answers.
EAB-P changed 55--64 answers.

The bias formula applies on rows $a\ge\tend$ (implementation name \texttt{ans1}).
A previous default applied the bias only on decode rows (\texttt{ans2}) and was labeled as the paper form.
That label is withdrawn.
The gate in those runs read $U_{\mathrm{post}}$, not the retrieval-time uncertainty named in the text.
Appendix~\ref{app:dca} records the local algebra of this retired bias.
It is not an accuracy result.

\subsection{Failure Cases and Format Sensitivity}
\label{sec:cases}
Failure counts for missing evidence, wrong modality, and format mismatch are \ph{}.
No case is filled in, because the v2 runs have not been executed.
Format sensitivity uses the single frozen prompt from Section~\ref{sec:setup}.
A comparison across prompt variants, if it is run, is a development-split description and is \ph{} here.
It is not a search over prompts on the test split.

%% ===============================================================
\section{Conclusion}
We describe \method{} v2, a per-question choice among direct answering, static text, static image, static mixed evidence, and reranked static evidence.
v2-Rule fits no weights.
v2-Cal is one logistic regression with seven parameters.
Fixed-sequence Learn-then-Test at $\alpha=\delta=0.10$ supplies the degradation-rate claim in Section~\ref{sec:risk}, which does not say that the method is never worse.
The primary comparison uses the strongest single action on the calibration split as the denominator.
The measurement, the main table, and the ablations have not been run, and their numbers are \ph{}.
The attention-steering subsection quotes the 13~October~2026 root-cause report: attention moves, and answers do not.
Appendix~\ref{app:limitations} states the limitations.
